\documentclass[12pt]{article}
\usepackage[utf8]{inputenc}
\usepackage[T1]{fontenc}
\usepackage{hyperref}
\title{Yiddish Language Resources for K-12 Curriculum Development}
\author{Compiled by Claude Sonnet 5.5 (claude-sonnet-5-5) for Andr\'as Kornai}
\date{}
\begin{document}
\maketitle

\section{Introduction}
Search found supplementary Yiddish culture materials for K-12 classes, mostly in English. It found no complete Yiddish-language K-12 curriculum. The Yiddish Book Center's own commentary notes that few schools still have Yiddish on the curriculum.

\section{Primary Educational Resources}
\begin{itemize}
\item KlezCalifornia Tam program: 26 free K-12 lesson plans on klezmer, Yiddish literature and food, tested in 25 schools. \url{https://ingeveb.org/pedagogy/tam-tastes-of-yiddish-culture-for-kids-teens}
\item Yiddish Book Center resource kits for teachers: \url{https://www.yiddishbookcenter.org/educational-programs/resources-teachers/resource-kits-teachers}
\item Yiddish Book Center programs and resources for teachers: \url{https://www.yiddishbookcenter.org/educational-programs/programs-and-resources-teachers}
\end{itemize}

\section{Linguistic and Reference Materials}
\begin{itemize}
\item Uriel Weinreich's dictionary (1968) and \emph{College Yiddish} (1949), as discussed by YIVO: \url{https://archives.yivo.org/When-Uriel-Weinreichs-Dictionary-Was-New}
\end{itemize}

\section{Religious and Community Resources}
\begin{itemize}
\item Workmen's Circle Boston Yiddish courses (2014 announcement): \url{https://mailman.yale.edu/pipermail/mendele/2014-September/000900.html}
\item Weekly Reader: The History of Yiddish Education (Yiddish Book Center): \url{https://www.yiddishbookcenter.org/language-literature-culture/weekly-reader/weekly-reader-history-yiddish-education}
\end{itemize}

\section{Digital and Audio Resources}
\begin{itemize}
\item YIVO-Bard summer program (Uriel Weinreich Program): \url{https://summerprogram.yivo.org/from-the-pages-of-yedies-38}
\end{itemize}

\section{Curriculum Development Context}
Existing materials are largely English-language culture lessons and adult or university language courses. A native-language K-12 curriculum would need new material. Sources suggest that Jewish day schools and Workmen's Circle schools are the places to look next, but this search did not confirm their current offerings.
\end{document}